\chapter{Objektorientierte Programmierung}\label{ch:K07_Klassen}
\section{Klassen}
Bisher haben wir mit einfachen Variablen (integer, string, boolean) und Datenstrukturen (Listen, Dictionaries) gearbeitet. Letztlich sind alle Python-Variablen Objekte, welche bestimmte Eigenschaften und Methoden besitzen können. Methoden haben wir bereits bei Listen kennengelernt, wo wir beispielsweise die Methode \lstinline|.append()| verwendet haben, um ein Element an eine Liste anzuhängen. Diese Methoden sind spezifisch für den Datentyp der Liste und ermöglichen es uns, auf einfache Weise mit den Daten zu interagieren.

Python erlaubt es uns, eigene Datentypen zu definieren, die sowohl Daten als auch Funktionen, also Methoden, enthalten können. Dies geschieht durch die Definition von \textbf{Klassen}, welche ein zentrales Konzept innerhalb des Paradigmas ``\gls{OOP}'' sind. Klassen ermöglichen es uns, komplexe Objekte zu modellieren, die sowohl Eigenschaften (Attribute) als auch Verhalten (Methoden) besitzen.

\begin{mydefinition}
Eine \textbf{Klasse} ist eine Vorlage für Objekte (Variablen), die Daten (Attribute) und Funktionen (Methoden) zusammenfasst. Sie definiert die Eigenschaften und das Verhalten einer Gruppe von Objekten. Ein \textbf{Objekt} ist eine konkrete Instanz (=eine konkrete Kopie) einer Klasse.

\textbf{Analogie}:
\begin{enumerate}
	\item Eine Klasse ist meist wie ein Bauplan für ein Haus. Der Bauplan beschreibt, wie das Haus aussehen soll (Attribute, z.B. Grösse, Farbe etc.) und welche Funktionen es hat (Methoden, etwa das Haus umbauen, abreissen etc.). Ein konkretes Haus, das nach diesem Plan gebaut wurde, ist eine Instanz dieser Klasse.
	\item Stellen Sie sich eine Klasse wie die Vorlage für einen Menschen  vor. Die Klasse beschreibt allgemeine Eigenschaften und Fähigkeiten des Menschen (z.B. Name, Körpergrösse, Interessensgebiete, etc.). Ein konkretes Objekt wäre dann ``Albert Einstein'', mit spezifischen Attributen wie \lstinline|name = "Albert Einstein"| und \lstinline|interessen = "Physik"|. So können Sie aus der Vorlage (Klasse) beliebig viele Menschen-Objekte mit individuellen Eigenschaften erzeugen.
\end{enumerate}
\end{mydefinition}

\begin{myexample}[Klassen definieren]
Eine Klasse wird mit dem Schlüsselwort \lstinline|class| definiert. Im folgenden Beispiel erstellen wir eine Klasse \lstinline|Person| mit den Attributen \lstinline|name| und \lstinline|alter| sowie einer Methode \lstinline|geburtstag_feiern|.

\lstinputlisting{Grundlagen_Info/00_Programmieren/Skript/Code/K07_Klassen/person.py}
\end{myexample}

In diesem Beispiel sehen wir mehrere wichtige Konzepte:

\begin{itemize}
	\item \lstinline|__init__| ist eine spezielle Methode (Konstruktor), die beim Erstellen eines Objekts aufgerufen wird. Hier werden die im Aufruf mitgegebenen Parameter \lstinline|name| und \lstinline|alter| den Attributen des Objekts, \lstinline|self.name| und \lstinline|self.alter|, zugewiesen.
	\item \lstinline|self| bezieht sich auf die aktuelle Instanz des Objekts und muss der erste Parameter jeder Methode sein.
	\item Objektattribute werden mit \lstinline|self.attributname| definiert und zugegriffen.
	\item Methoden werden mit \lstinline|objekt.methode()| aufgerufen.
\end{itemize}

\begin{myexample}[Bankkonto-Klasse]
Ein praktisches Beispiel ist eine Klasse für Bankkonten:

\lstinputlisting{Grundlagen_Info/00_Programmieren/Skript/Code/K07_Klassen/bankkonto.py}
\end{myexample}

\begin{myexercise}
	Erstellen Sie eine Klasse \lstinline|Buch| mit den Attributen \lstinline|titel|, \lstinline|autor| und \lstinline|seitenzahl|. Implementieren Sie eine Methode \lstinline|info()|, die eine Zusammenfassung der Buchinformationen ausgibt. Erstellen Sie zwei Buchobjekte und rufen Sie die \lstinline|info()|-Methode auf.

	\begin{myanswer}
		\lstinputlisting{Grundlagen_Info/00_Programmieren/Skript/Code/K07_Klassen/buch.py}
	\end{myanswer}
\end{myexercise}

\begin{myexercise}
	Erweitern Sie die \lstinline|Bankkonto|-Klasse um eine Methode \lstinline|ueberweisen(zielkonto, betrag)|, die Geld von einem Konto auf ein anderes überweist. Die Methode soll prüfen, ob genügend Guthaben vorhanden ist, und dann den Betrag vom aktuellen Konto abziehen und zum Zielkonto hinzufügen.

	\begin{myanswer}
		\lstinputlisting{Grundlagen_Info/00_Programmieren/Skript/Code/K07_Klassen/bankkonto_ueberweisung.py}
	\end{myanswer}
\end{myexercise}

\section{Vordefinierte Klassen in Python}

Wie bereits erwähnt ist jegliche Variable in Python ein Objekt, das zu einer bestimmten Klasse gehört. Python bietet eine Vielzahl von vordefinierten Klassen, welche wir in den vorausgehenden Kapiteln bereits angeschaut haben, beispielsweise Datentypen wie \lstinline|int|, \lstinline|float|, \lstinline|str| (String), \lstinline|list| (Liste) oder \lstinline|dict| (Dictionary), sowie weitere Klassen. Diese Klassen haben ihre eigenen Methoden, die spezifische Operationen auf den Objekten dieser Klassen ermöglichen.

\begin{myremark}[Vordefinierte Klassen in Python]
	Eine Auflistung aller vordefinierten Klassen in Python finden Sie in der offiziellen Dokumentation unter \url{https://docs.python.org/3/library/stdtypes.html}. Dort sind alle Datentypen und deren Methoden beschrieben, die in Python verfügbar sind. Alternativ kann folgender Code ausgeführt werden, um die Namen aller vordefinierten Klassen in Python zu erhalten:
	\lstinputlisting{Grundlagen_Info/00_Programmieren/Skript/Code/K07_Klassen/builtin_classes.py}
\end{myremark}

Weshalb kann man in Python eine Variable für gewisse vordefinierte Klassen erstellen, ohne explizit den Namen der Klasse anzugeben? Hier handelt es sich um etwas ``Python-Magie'', bei der Python den Typ der Variable automatisch erkennt und die entsprechende Klasse verwendet. Wenn wir beispielsweise eine Variable \lstinline|x = 5| erstellen, wird \lstinline|x| automatisch als Objekt der Klasse \lstinline|int| erstellt. Python kümmert sich im Hintergrund um die Zuweisung der richtigen Klasse, sodass wir uns nicht explizit darum kümmern müssen. Wir könnten jedoch die Klasse explizit angeben, indem wir beispielsweise \lstinline|x = int(5)| schreiben, was dasselbe Ergebnis liefert.

Wie wir gesehen haben, beinhalten gewisse Klassen bereits Methoden, wie beispielsweise \lstinline|.append()| für Listen oder \lstinline|.keys()| für Dictionaries. Diese Methoden sind spezifisch für die jeweilige Klasse und ermöglichen es uns, auf einfache Weise mit den Objekten dieser Klassen zu interagieren. Falls wir wissen möchten, welche Methoden eine bestimmte Klasse hat, können wir die Funktion \lstinline|dir()| verwenden, um eine Liste aller verfügbaren Methoden und Attribute zu erhalten. Zum Beispiel:

\begin{myexample}[Verfügbare Methoden einer Klasse] Folgendes Beispiel zeigt, wie wir die verfügbaren Methoden für eine Liste und einen String abfragen können. Dies kann sowohl für eine konkrete Instanz wie auch für den Klassennamen selbst erfolgen:
\lstinputlisting{Grundlagen_Info/00_Programmieren/Skript/Code/K07_Klassen/verfuegbare_methoden.py}
\end{myexample}

\section{Klassenmethoden und Attribute}

Neben den Instanzattributen (die zu jedem Objekt gehören) können Klassen auch Klassenattribute haben, die für alle Instanzen gleich sind.

\begin{myexample}[Klassenattribute] Folgendes Beispiel zeigt, wie Klassenattribute definiert und verwendet werden können, beispielsweise um SchülerInnen mit Attributen wie Name und Klasse zu erstellen:
\lstinputlisting{Grundlagen_Info/00_Programmieren/Skript/Code/K07_Klassen/klassenattribute.py}
\end{myexample}

\begin{myexample}[Instanzzähler]
Ein häufiger Anwendungsfall für Klassenattribute ist das Zählen von Instanzen:

\lstinputlisting{Grundlagen_Info/00_Programmieren/Skript/Code/K07_Klassen/instanzzaehler.py}
\end{myexample}

\begin{myexercise}
	Erstellen Sie eine Klasse \lstinline|Auto| mit den Instanzattributen \lstinline|marke|, \lstinline|modell| und \lstinline|baujahr|. Fügen Sie ein Klassenattribut \lstinline|anzahl_autos| hinzu, das die Gesamtzahl der erstellten Auto-Objekte zählt. Implementieren Sie eine Klassenmethode \lstinline|get_statistik()|, die die Anzahl der Autos zurückgibt.

	\begin{myanswer}
		\lstinputlisting{Grundlagen_Info/00_Programmieren/Skript/Code/K07_Klassen/auto.py}
	\end{myanswer}
\end{myexercise}

\section{Vererbung und Polymorphismus}
\subsection{Vererbung}

Ein mächtiges Konzept der Objektorientierung ist die Vererbung. Sie ermöglicht es, eine neue Klasse basierend auf einer vorhandenen Klasse zu erstellen und deren Eigenschaften und Methoden zu übernehmen.

\begin{mydefinition}
Bei der \textbf{Vererbung} erbt eine Kindklasse (Subklasse) Attribute und Methoden von einer Elternklasse (Superklasse). Die Kindklasse kann zusätzliche Attribute und Methoden haben oder vorhandene überschreiben. Die Vererbung geschieht mithilfe des Schlüsselworts \lstinline|class Kindklasse(Elternklasse):| sowie der Verwendung von \lstinline|super()|, um auf die Elternklasse zuzugreifen. Dabei wird zuerst eine Instanz der Elternklasse erstellt, bevor die Kindklasse ihre eigenen Attribute und Methoden hinzufügt oder überschreibt.
\end{mydefinition}

\begin{myexample}[Vererbung]
\lstinputlisting{Grundlagen_Info/00_Programmieren/Skript/Code/K07_Klassen/elektroauto.py}
\end{myexample}

In diesem Beispiel:
\begin{itemize}
	\item \lstinline|ElektroAuto| erbt von \lstinline|Fahrzeug| und erhält alle seine Attribute und Methoden.
	\item Mit \lstinline|super().__init__(...)| rufen wir den Konstruktor der Elternklasse auf.
	\item Die Methode \lstinline|fahren()| wird in der Kindklasse überschrieben, um die spezifische Funktionalität eines Elektroautos zu implementieren.
	\item Mit \lstinline|super().info()| greifen wir auf die Methode der Elternklasse zu.
\end{itemize}

\begin{myexercise}
	Erstellen Sie eine Elternklasse \lstinline|Person| mit den Attributen \lstinline|name| und \lstinline|alter| sowie einer Methode \lstinline|geburtstag_feiern()|. Erstellen Sie dann eine Kindklasse \lstinline|Schueler| mit einem zusätzlichen Attribut \lstinline|klasse| und einer überschriebenen Methode \lstinline|vorstellen()|, die auch die Klasse ausgibt.

	\begin{myanswer}
		\lstinputlisting{Grundlagen_Info/00_Programmieren/Skript/Code/K07_Klassen/schueler_child.py}
	\end{myanswer}
\end{myexercise}

\begin{myexercise}
	Erstellen Sie eine Basisklasse \lstinline|Bankkonto| wie im vorherigen Beispiel. Dann erstellen Sie eine Unterklasse \lstinline|Sparkonto|, die eine zusätzliche Methode \lstinline|zinsen_gutschreiben(zinssatz)| hat, die Zinsen basierend auf dem aktuellen Kontostand gutschreibt.

	\begin{myanswer}
		\lstinputlisting{Grundlagen_Info/00_Programmieren/Skript/Code/K07_Klassen/sparkonto.py}
	\end{myanswer}
\end{myexercise}

\subsection{Polymorphismus}
Unterschiedliche Klassen können Methoden mit denselben Namen haben, aber jede Klasse hat ihre eigene Implementierung. Dies wird als \textbf{Polymorphismus} bezeichnet. Das Wort Polymorphismus stammt aus dem Griechischen und bedeutet soviel wie "viele Formen", von gr. \textit{poly} = viel und \textit{morphos} = Form. Polymorphismus ermöglicht, dass verschiedene Klassen auf die gleiche Weise angesprochen werden können, obwohl sich die Methoden unterscheiden. Dies ist insbesondere bei Kindesklassen von Bedeutung, die die Methoden der Elternklasse überschreiben oder auf unterschiedliche Weise implementieren können.

\begin{myexample}[Polymorphismus] Folgendes Beispiel zeigt, wie verschiedene Klassen die gleiche Methode \lstinline|geraeusch()| implementieren können:
\lstinputlisting{Grundlagen_Info/00_Programmieren/Skript/Code/K07_Klassen/polymorphism.py}
\end{myexample}


\section{Praktisches Beispiel: Bibliothekssystem}

Als umfassenderes Beispiel implementieren wir ein einfaches Bibliothekssystem mit mehreren Klassen, die verschiedene Aspekte einer Bibliothek modellieren.

\begin{myexample}[Bibliothekssystem]Folgendes Beispiel zeigt, wie wir Klassen für Bücher, Autoren usw. für eine eine Bibliothek erstellen können:
\lstinputlisting{Grundlagen_Info/00_Programmieren/Skript/Code/K07_Klassen/library_system.py}
\end{myexample}

\begin{myexercise}
	Erweitern Sie das Bibliothekssystem um eine neue Klasse \lstinline|Zeitschrift|, die von \lstinline|Buch| erbt, aber zusätzlich eine \lstinline|ausgabe| (z.B. "Mai 2024") hat. Überschreiben Sie die \lstinline|info()|-Methode entsprechend, und fügen Sie mindestens eine Zeitschrift zur Bibliothek hinzu.

	\begin{myanswer}
		\lstinputlisting{Grundlagen_Info/00_Programmieren/Skript/Code/K07_Klassen/zeitschrift.py}
	\end{myanswer}
\end{myexercise}

\begin{mychallenge}
	Entwickeln Sie ein vollständiges Lagerverwaltungssystem mit folgenden Klassen:
	\begin{itemize}
		\item \lstinline|Produkt| (Basisklasse mit Name, Artikelnummer, Preis)
		\item \lstinline|Elektronikprodukt| (erbt von Produkt, zusätzlich mit Garantiedauer)
		\item \lstinline|Lebensmittel| (erbt von Produkt, zusätzlich mit Haltbarkeitsdatum)
		\item \lstinline|Lager| (verwaltet Produkte, mit Methoden zum Hinzufügen, Entfernen, Suchen und Bestandsanzeige)
		\item \lstinline|Bestellung| (enthält Produkte und Mengen, berechnet Gesamtpreis)
	\end{itemize}

	Implementieren Sie auch eine Methode, die abgelaufene Lebensmittel identifiziert und aus dem Lager entfernt.
\end{mychallenge}



\section{Zusammenfassung}
Die objektorientierte Programmierung mit Klassen ist ein mächtiges Werkzeug zur Strukturierung von Code und Modellierung realer Objekte. Wichtige Konzepte sind:

\begin{itemize}
	\item Klassen definieren Vorlagen für Objekte mit Attributen und Methoden
	\item Objekte sind konkrete Instanzen von Klassen
	\item Methoden ermöglichen es, auf Attribute zuzugreifen und diese zu verändern.
	\item Vererbung ermöglicht die Wiederverwendung und Erweiterung von Code in Form von Kindklassen, die ihre Attribute von Elternklassen erben und gegebenenfalls überschreiben.
	\item Polymorphismus ermöglicht es, Methoden mit denselben Namen in verschiedenen Klassen zu verwenden, wobei jede Klasse ihre eigene, unabhängige Implementierung hat.
\end{itemize}

Durch die Verwendung von Klassen können komplexe Systeme modelliert und implementiert werden, wodurch der Code besser strukturiert, lesbarer und wartbarer wird. Insbesondere in Games sind Klassen ein zentrales Konzept, um verschiedene Spielobjekte (wie Spieler, Gegner, Items) zu modellieren und deren Verhalten zu steuern. In der Praxis werden Klassen häufig in Kombination mit anderen Konzepten wie Vererbung und Polymorphismus verwendet, um flexible und erweiterbare Softwarearchitekturen zu schaffen.
